\documentstyle[%


righttag%


,11pt%


,amssymb%


,russian%               remove in Eng.


,izvru%                 Set izven% in Eng.


% ,izvru-ks             for brief communications


]{amsart}





% Touched 04.10 Verify enumeration of eqs.


\renewcommand{\d}{{\mathrm{d}}}


\newcommand{\ord}{\operatorname{ord}}


\newcommand{\tts}[1]{}





%\hoffset=-15mm%                  For Eng.


\setcounter{page}{19}


\god{2000}


\nomer{2~(453)} %                        Remove in Eng.


%\nomer{Vol.\,44, No.\,2}%               Set in Eng.


%\str{\pageref{begin}--\pageref{end}}%  Set in Eng.


\udk{517.544}





\author{\it И.А.\,Бикчантаев}


% NB:   ^^ remove in Eng.


\title{Задача Римана на ультрагиперэллиптической поверхности}


%\thanks{Работа выполнена при поддержке фонда НИОКР Республики


%Татарстан (грант 09-12/1999).}





\date{}


%\pagestyle{myheadings}%         Set in Eng.


\pagestyle{plain} %              Remove in Eng.


\begin{document}


%\thispagestyle{plain}%          Set in Eng.


\maketitle


\label{begin}





В работах [1], [2]


было получено в квадратурах решение краевой задачи Римана на


компактной римановой поверхности. В работах [3] и [4] на некомпактной


римановой поверхности была построена ее н\"етерова теория и вычислены


дефектные числа. На произвольной открытой римановой поверхности эта


задача была решена в явном виде лишь в случаях, когда коэффициент


задачи равен 1 [3] или $-1$ [5].


В данной работе получены условия разрешимости и дано


явное решение краевой задачи Римана на ультрагиперэллиптической


поверхности $R$ в случае произвольного кусочно-гладкого контура


$\Gamma$. При этом решение (в общем случае) выражается через


интегралы, ядрами которых служат аналоги ядра Коши на некоторой


вспомогательной гиперэллиптической поверхности. В том случае, когда


контур $\Gamma$ не разбивает поверхность $R$, оказалось, что можно


обойтись более элементарными средствами, чем при решении задачи Римана


на компактной поверхности (в частности, гиперэллиптической). А именно,


не используется проблема обращения Якоби (или какой-либо ее аналог),


вместо аналогов ядра Коши на римановой поверхности используется


обычное ядро Коши на комплексной плоскости.


Отметим, что на такой поверхности


условие $\Lambda_0$-поведения искомой функции (которое использовалось


в [3]--[5]) означает ее ограниченность в окрестности


идеальной границы.





\section{Предварительные сведения и обозначения}





Пусть $(R,z)$ есть безграничная двулистная накрывающая комплексной


плоскости $\Bbb C$, точки ветвления которой сгущаются к единственной


точке идеальной границы, лежащей над $z=\infty$. Накрывающая $(R,z)$ может


быть реализована как риманова поверхность функции $u(z)$, определяемой


уравнением $u^2=P(z)$, где $P(z)$ --- целая функция с бесконечным


числом простых


нулей. Риманова поверхность $R$, допускающая такую реализацию, называется


ультрагиперэллиптической. Через $j_R$ обозначим отличное от тождественного


преобразование наложения


накрывающей $(R,z)$, т.\,е. конформный автоморфизм поверхности $R$,


удовлетворяющий соотношению $z(j_R(q))=z(q)$ для всех $q\in R$. Пусть


$\Gamma$ --- кусочно-гладкая линия на $R$ и $\zeta (q)$ --- голоморфная


функция в окрестности $\Gamma$


такая, что $d\zeta$ не имеет нулей на $\Gamma$. Тогда $\zeta$ является


локальной униформизирующей в окрестности любой точки $\tau \in \Gamma$.


Если $\Gamma$ не содержит точек ветвления накрывающей $(R,z)$, то можно


взять $\zeta=z$. Если $\tau\in\Gamma$ --- точка ветвления, то в ее


окрестности $z$ и $\zeta$ связаны соотношением вида $z-z(\tau)=(\zeta


-\zeta(\tau))^2a(\zeta)$, где $a(\zeta)$ --- функция, голоморфная в точке


$\zeta =\zeta (\tau)$ и $a(\zeta (\tau))\ne 0$.





Пусть $\alpha:[0, 1] \to R $ --- путь на $R$. Его длиной


назовем длину плоского пути $z\circ\alpha:[0,1]\to\Bbb C$,


т.\,е. $\int\limits_0^1|\mathrm{d}(z\circ\alpha)|$. Расстоянием $\rho(p,q)$


между точками $p$ и $q$ на $R$ назовем точную нижнюю грань длин всех


путей, соединяющих $p$ и $q$. Используя локальную однолистность


функции $\zeta(q)$, нетрудно показать, что найдется число $\epsilon


>0$ такое, что из неравенств $\rho(t_1,t_2)\le\epsilon$, $t_1\ne t_2$


вытекает неравенство $\zeta(t_1)\ne\zeta(t_2)$, $t_1,t_2\in\Gamma$.





Обозначим через $T$ множество всех узлов


контура $\Gamma$. На $T$ определим


действительнозначную функцию $\lambda=\lambda(\tau)$. По аналогии с [6]


введем пространство $H_{\mu,\lambda}(\Gamma,T)$, состоящее из функций


на~$\Gamma$, $H_\mu$-непрерывных вне любой окрестности множества $T$ и


ведущих себя вблизи $T$ как весовая функция


$$


\rho_\lambda (t)=\prod_{\tau\in T}|\zeta (t)-\zeta(\tau)|^{\lambda(\tau)}.


$$


Пространство $H_{\mu,\lambda}(\Gamma, T)$ банахово относительно нормы


$$


\|\varphi \|_{\mu,\lambda}=


\sup_{t\in \Gamma}|\varphi (t)\rho_{-\lambda}(t)|+


\{\rho_{\mu -\lambda}\varphi\}_\mu,


$$


где


$$


\{\varphi\}_{\mu}=


\sup_{t_1,t_2\in\Gamma, \ \rho (t_1, t_2)<\epsilon,\ t_1 \ne t_2}


\frac{|\varphi (t_1)-\varphi (t_2)|}


{|\zeta (t_1)-\zeta (t_2)|^\mu}\,.


$$


Выберем $\delta>0$ столь малым, чтобы множества


$$


\Gamma_\tau =\Gamma\cap\{q\in R: 0<\rho (q, \tau)\le\delta\},


\ \; \tau\in T,$$


попарно не пересекались и распадались на компоненты $\Gamma_{\tau, i}$,


$1\le i\le n_\tau$, $\tau \in T$. Введем класс


$H_{\mu,(\lambda)}(\Gamma, T)$\ функций $\varphi (t)$, которые


$H_\mu$-непрерывны на $\Gamma$ вне любой окрестности множества $T$ и на


каждом $\Gamma_\tau$, $\tau \in T$, представимы в виде


$$


\varphi (t)=p_{\tau, i}(t)+\varphi_\tau (t), \ \; t \in \Gamma_{\tau, i},


$$


где $p_{\tau, i}$ --- многочлен степени меньше, чем $\lambda(\tau)$


(многочлен отрицательной степени условимся считать нулем),


$\varphi_\tau\in H_{\mu,\lambda(\tau)}(\Gamma_\tau, \tau)$. Пространство


$H_{\mu,(\lambda)}(\Gamma,T)$ банахово относительно нормы


$$


\|\varphi \|_{\mu,(\lambda)}=\sum_{\tau \in T}


\bigg(\,\sum_{i=1}^{n_\tau}\sup_{\Gamma_\tau}|p_{\tau,i}|


+\|\varphi_\tau\|_{H_{\mu,\lambda(\tau)}(\Gamma_\tau,\tau)}\bigg)+


\|\varphi\|_{H_{\mu}(\Gamma')},


$$


где $\Gamma'=\Gamma\cap\{q:\rho (q, T)\ge \delta /2\}$,


$\rho (q,T)=\min_{\tau\in T}\rho (q, \tau)$ (ср. [6]).





\section{Случай неразбивающего контура}





1. Пусть $\Gamma$ --- кусочно-гладкий


контур на $R$ такой, что множество $R\setminus \Gamma$ связно. Будем


предполагать, что $z(\Gamma)$ тоже есть кусочно-гладкий контур на


$\Bbb C$, причем $\Gamma$ и $z(\Gamma)$ не имеют точек возврата.


Ориентацию на $\Gamma$ выберем так, чтобы на линейных участках


контура, имеющих одинаковые проекции на $z$-плоскость $\Bbb C$,


ориентация контура была согласована относительно преобразования


наложения $j$. Тогда ориентацию $z(\Gamma)$ можно определить как


индуцированную отображением $z:\Gamma \to z(\Gamma)$.





Зададим дивизор $D$, носитель которого лежит в $R\setminus \Gamma$, и


функции $G \in H_{\mu, (\mu)}(\Gamma, T)$ и


$g \in H_{\mu, \lambda}(\Gamma, T)$, $0<\mu < 1$, $-1<\lambda <0$,


причем $G(t)\ne 0$, $t\in \Gamma$. Рассмотрим краевую задачу Римана в


следующей постановке.





Найти кусочно-мероморфную функцию $F$ на $R$ с линией скачков


$\Gamma$, кратную дивизору $1/D$ и ограниченную в окрестности идеальной


границы поверхности $R$, предельные значения которой на $\Gamma$


принадлежат классу $H_{\mu, \lambda}(\Gamma, T)$, $0<\mu <1$,


$-1<\lambda<0$, и удовлетворяют соотношению


\begin{equation}


F^+(t)=G(t)F^-(t)+g(t),\quad t\in \Gamma.


\end{equation}





2. Через $q(z)$ будем обозначать точку на $R$ такую, что


$z(q(z))=z$ (поднятие точки $z\in\Bbb C$ на накрывающую $(R,z))$.


Если $F$ --- решение задачи (1), то функция


$(F(j(q(z)))-F(q(z)))^2$ голоморфно продолжима в точку $z=\infty$,


являющуюся точкой сгущения ее нулей. Следовательно, $F\circ j=F$ (ср.


[7]) и функция $f(z)=F(q(z))$ однозначна в $\Bbb C\setminus


z(\Gamma)$ и аналитически продолжима в $\Bbb C\setminus\gamma$,


где $\gamma$ есть множество точек на $z(\Gamma)$, имеющих два


прообраза при отображении $z: \Gamma \to z(\Gamma)$ (точки ветвления


накрывающей $(R, z)$ при этом считаются дважды).





3. Множество $M$ на римановой поверхности или на плоскости называется


$AB$-устранимым, если для некоторой окрестности $U$ множества $M$ любая


аналитическая и ограниченная в $U\setminus M$ функция аналитически


продолжима в $U$. Определим в $\Bbb C$ дивизор $\Delta$, полагая


$\ord_{z(q)}\Delta =\min (\ord_qD,\ \ord_{j(q)}D)$


при $j(q)\ne q$ и $\ord_{z(q)}\Delta =[\frac12\ord_qD]$ при


$j(q)=q$, где $[\ ]$ означает целую часть числа. Если множество $\gamma$


является $AB$-устранимым, то функция $f$ аналитически продолжима на всю


комплексную плоскость $\Bbb C$ и является рациональной функцией,


кратной дивизору $1/\Delta$. Из условия (1) вытекает, что в случае


разрешимости задачи Римана функции $G$ и $g$ удовлетворяют соотношению


\begin{equation}


g(t)=(1-G(t))f(z(t)),\quad t\in \Gamma,


\end{equation}


где $f$ --- рациональная функция, кратная дивизору $1/\Delta$. Если это


соотношение выполняется, то функция $F(q)=f(z(q))$ является решением краевой


задачи Римана (1). Отсюда вытекает





\begin{thm}


Пусть $\gamma$ --- $AB$-устранимое множество.


Тогда для разрешимости краевой задачи Римана $(1)$ необходимо и достаточно,


чтобы функции $G$ и $g$ были связаны соотношением $(2)$, где $f$ ---


рациональная функция, кратная дивизору $1/\Delta$.


\end{thm}





\begin{cor}


При $\ord\Delta<0$ задача (1) имеет решение (равное


нулю) только при $g=0$.


\end{cor}





\begin{cor}


Если $G=1$, то для разрешимости задачи Римана (1)


необходимо и достаточно, чтобы $g=0$. При этом любое решение задачи (1)


имеет вид $F(q)=f(z(q))$, где $f$ --- рациональная функция, кратная дивизору


$1/\Delta$. Число линейно независимых решений равно


$\max (0,\ \ord\Delta+1)$.


\end{cor}





\begin{cor}


Если $G(t)\not\equiv 1$, то задача (1) не может иметь более


одного решения.


\end{cor}





Действительно, при $g=0$ условие (2) может выполняться только при


$f=0$. Поэтому соответствующая однородная задача (1) имеет лишь нулевое


решение.





4. Предположим теперь, что множество $\gamma$ не является


$AB$-устранимым; при этом его линейная мера будет положительной. Обозначим


через $\Gamma_1$ кривую на $R$, гомеоморфную $z(\Gamma)$ относительно


отображения $z:\Gamma_1\to z(\Gamma)$. Тогда $\Gamma_2:=j(\Gamma_1)$


тоже гомеоморфна $z(\Gamma)$ относительно отображения $z$ и


$\Gamma_1\cup\Gamma_2=z^{-1}(z(\Gamma))$. Через $\rho_k:z(\Gamma)\to


\Gamma_k$, $k=1,\,2$, обозначим гомеоморфизм $z(\Gamma)$ на $\Gamma_k$


такой, что $z(\rho_k(\xi))=\xi$ при $\xi \in z(\Gamma)$. Ориентацию на


$\Gamma_k$ выберем таким образом, чтобы индуцированная ею при


проектировании $z:R\to\Bbb C$ ориентация на $z(\Gamma)$ совпадала с


уже выбранной в п.\,1.





Доопределим $G$ и $g$ на $\Gamma_1 \cup \Gamma_2$, полагая $G=1,


g=0$ в точках, не принадлежащих $\Gamma$. Тогда функция $f$ на $z(\Gamma)$


удовлетворяет условиям


\begin{equation}


f^+(\xi)=G(\rho_k(\xi))f^-(\xi)+g(\rho_k(\xi)),\quad \xi \in z(\Gamma), 


\quad \Delta^{-1}|(f),\quad k=1,\ 2.\tag{$3_k$}


\end{equation}


\addtocounter{equation}{1}





\noindent Здесь функции $G(\rho_k(\xi))$ принадлежат классу


$H_{\mu, (\mu')}(z(\Gamma), z(T))$, где


$$


\mu'(z(\tau))=


\begin{cases}


\mu, & \tau \in T,\ \; j(\tau)\ne \tau; \\


\frac12\mu, & \tau \in T,\ \; j(\tau)=\tau.


\end{cases}


$$


Функции $g(\rho_k(\xi))$ принадлежат классу $H_{\mu,\nu}(z(\Gamma), z(T))$,


где


$$


\nu (z(\tau))=


\begin{cases}


\min (\lambda(\tau), \lambda (j(\tau))), &


\tau,\, j(\tau)\in T,\ \; j(\tau)\ne \tau; \\


\lambda (\tau), & \tau \in T,\ \; j(\tau)\notin T; \\


\frac12\lambda(\tau), & j(\tau)=\tau\in T.


\end{cases}


$$


Решение задачи $(3_k)$ будем отыскивать в классе функций, предельные значения


которых на $z(\Gamma)$ принадлежат классу $H_{\mu, \nu}(z(\Gamma), z(T))$.





Таким образом, функция $f$ является одновременно решением двух


краевых задач Римана на контуре $z(\Gamma)$. Из связности $R\setminus\Gamma$


следует связность множества $\Bbb C\setminus \gamma$. В точки


множества\linebreak


$z(\Gamma)\setminus \gamma$ функция $f$ аналитически продолжима и,


следовательно, является аналитической в области $\Bbb C\setminus \gamma$


за исключением, возможно, конечного числа полюсов. Поэтому для совпадения


решений краевых задач $(3_1)$ и $(3_2)$ достаточно потребовать их


совпадения в окрестности некоторой точки


$z_0\in\Bbb C\setminus z(\Gamma)$.





5. Если $G\circ j=G$ на $\Gamma_1\cup\Gamma_2$, то, поскольку


$F\circ j=F,$ для разрешимости задачи (1) необходимо, чтобы $g\circ j=g$


на $\Gamma_1\cup\Gamma_2$. Тогда задачи $(3_1)$ и $(3_2)$


совпадают. Если $f$ --- решение задачи $(3_k)$, то


$F(q)=f(z(q))$ будет решением задачи (1).





Обозначим через $X(z)$ каноническую функцию задачи $(3_k)$,


предельные значения которой на $\Gamma$ принадлежат классу


$H_{\mu, \nu}(z(\Gamma), z(T))$ и которая имеет максимально возможный


порядок $\kappa$ на бесконечности. Тогда при $\kappa+\ord\Delta\ge -1$


задача $(3_k)$ разрешима при любой функции $g\in


H_{\mu, \lambda}(\Gamma,T)$, $0<\mu <1$, $-1<\lambda <0$,


и ее общее решение имеет вид


\begin{equation}


f(z)=\frac{X(z)}{2\pi i}\int_{z(\Gamma)}\frac{g(\rho_k(\xi))d\xi}


{X^+(\xi)(\xi -z)} +X(z)\delta (z),


\end{equation}


где $\delta$ --- произвольная рациональная функция, кратная дивизору


$\Delta^{-1}\infty^{-\kappa}$.





При $\kappa + \ord\Delta <-1$ необходимые и достаточные условия


разрешимости задачи $(3_k)$, а следовательно, и задачи (1) имеют вид


\begin{equation}


\int_{z(\Gamma)}\frac{g(\rho_k(\xi))\omega_j(\xi)d\xi}{X^+(\xi)}=0,


\quad j=1,\dotsc, -\kappa -\ord\Delta -1,


\end{equation}


где $\omega_j$ --- базис пространства рациональных функций, кратных дивизору


$\Delta\infty^{\kappa +2}$. При выполнении условий (5) задача $(3_k)$


имеет единственное решение вида (4), где $\delta =0$.





Условия (5) можно переписать в виде


\begin{equation}


\int_{\Gamma\cap\Gamma_k}g\theta_j=0,\quad j=1,


\dotsc,-\kappa -\ord\Delta -1,


\end{equation}


где $\theta_j(t)=(X^+(z(t)))^{-1}\omega_j(z(t))dz(t)$. Общее решение задачи


(1) имеет вид


\begin{equation}


F(q)=\frac{X(z(q))}{2\pi i}\int_{\Gamma\cap\Gamma_k}\frac{g(t)dz(t)}


{X^+(z(t))(z(t)-z(q))}+X(z(q))\delta(z(q)).


\end{equation}


В формулах (6) и (7) $k$ может принимать любое из значений 1 или 2.


Таким образом, доказана 





\begin{thm}


Пусть $\gamma$ имеет положительную линейную меру и


коэффициент $G$ задачи $(1)$ удовлетворяет соотношению $G\circ j=G$ на


$\Gamma_1\cup\Gamma_2$. Тогда для разрешимости задачи $(1)$ необходимо и


достаточно, чтобы функция $g$ удовлетворяла условиям $g\circ j=g$ на


$\Gamma_1\cup\Gamma_2$ и $(6)$. При их выполнении общее решение задачи $(1)$


имеет вид $(7)$, где $\delta$ --- произвольная рациональная функция, кратная


дивизору $\Delta^{-1}\infty^{-\kappa}$. Однородная $(g=0)$ задача $(1)$ имеет


$l=\max(0,\ \ord\Delta +\kappa +1)$ линейно независимых решений.


\end{thm}





6. Пусть $\gamma$ такое же, как в предыдущем пункте. Здесь будем


предполагать, что $G\circ j\ne G$. Ясно, что при этом однородная


$(g=0)$ задача Римана (1) имеет лишь нулевое решение, а неоднородная может


иметь не более одного решения. В рассматриваемом случае функция $f$ является


решением одновременно двух различных краевых задач Римана $(3_1)$ и


$(3_2)$.


Обозначим через $X_k(z)$ каноническую функцию (того же класса, что и в


п.\,5) задачи $(3_k)$, $\kappa_k=\ord_\infty X_k(z)$.


При $\kappa_k +\ord\Delta\ge -1$ задача $(3_k)$


безусловно разрешима, и ее общее решение имеет вид


\begin{equation}


f_k(z)=\frac{X_k(z)}{2\pi


i}\int_{z(\Gamma)}\frac{g(\rho_k(\xi))d\xi}{X^+_k(\xi)(\xi -z)}+


X_k(z)\delta_k(z),


\end{equation}


где $\delta_k(z)$ --- произвольная рациональная функция, кратная дивизору


$\Delta^{-1}\infty^{-\kappa_k}$. При $\kappa_k +\ord\Delta <-1$


для разрешимости задачи $(3_k)$ необходимо и достаточно


выполнения условий


\begin{equation}


\int_{z(\Gamma)}\frac{g(\rho_k(\xi))}{X_k^+(\xi)}\omega_{kj}(\xi)d\xi


=0,\quad j=1,\dotsc,-\kappa_k-\ord\Delta -1,


\end{equation}


где $\omega_{kj}$ --- базис пространства рациональных функций, кратных


дивизору $\Delta\infty^{\kappa_k+2}$. Полагая


$$


\theta_{kj}(t)=\frac{\omega_{kj}(z(t))dz(t)}{X^+_k(z(t))},


$$


условие (9) перепишем в виде


\begin{equation}


\int_{\Gamma\cap\Gamma_k}g\theta_{kj}=0,\quad j=


1,\dotsc,-\kappa_k-\ord\Delta -1,\quad k=1,\,2.


\end{equation}


При выполнении условий (10) задача $(3_k)$ имеет единственное


решение вида (8), где $\delta_k=0$. Для того чтобы функции $f_k$


определяли решение исходной задачи (1), должно выполняться равенство


$f_1=f_2$. В силу сказанного в п.\,4 достаточно потребовать выполнения


этого равенства в окрестности некоторой точки $z_0\in


\Bbb C\setminus z(\Gamma)$.





Пусть $z_0$ --- точка из $\Bbb C\setminus z(\Gamma)$, в


которой функции $X_k$ и $f_k$ голоморфны. Используем


ряд Тейлора в окрестности точки $z_0$


$$


\frac{X_k(z)}{2\pi iX^+_k(\eta)(\eta -z)}=\sum_{j=0}^\infty c_{kj}(\eta)


(z-z_0)^j.


$$


Пусть $\delta_{kj}$, $j=1,\dotsc,\kappa_k +\ord\Delta +1$, ---


базис пространства рациональных функций, кратных дивизору


$\Delta^{-1}\infty^{-\kappa_k}$, $k=1,\,2$. Тогда


$$


X_k(z)\delta_k(z)=\sum_{n=1}^{\kappa_k +\ord\Delta +1}


a_{kn}X_k(z)\delta_{kn}(z),


$$


где $a_{kn}$ --- комплексные числа. Запишем для функции $X_k\delta_{kn}$


ряд Тейлора


$$


X_k(z)\delta_{kn}(z)=\sum_{j=0}^{\infty}a_{knj}(z-z_0)^j.


$$


Сравнивая коэффициенты тейлоровских разложений функций $f_1$ и $f_2$,


получим соотношения, эквивалентные равенству $f_1=f_2$,


\begin{equation}


\sum_{n=1}^{\kappa_1+\ord\Delta +1}a_{1n}a_{1nj}-


\sum_{n=1}^{\kappa_2+\ord\Delta +1}a_{2n}a_{2nj}=


-\int_{z(\Gamma)}(g(\rho_1(\eta))c_{1j}(\eta)-


g(\rho_2(\eta))c_{2j}(\eta))d\eta,\ \; j=0,1,2,\dotsc


\end{equation}


Обозначим через $A$ матрицу системы (11) и положим


\begin{gather*}


\alpha_j(t)=


\begin{cases}


-c_{1j}(z(t))dz(t), & t\in\Gamma\cap\Gamma_1;\\


\phantom{-}c_{2j}(z(t))dz(t), & t\in\Gamma\cap\Gamma_2,


\end{cases}


\\


\alpha (t)=(\alpha_1(t),\ \alpha_2(t),\dotsc)^t, \\


%


a=(a_{11},\dotsc,a_{1,\, \kappa_1+\ord\Delta +1},


\ a_{21},\dotsc,a_{2,\, \kappa_2+\ord\Delta +1})^t.


\end{gather*}


Тогда система (11) запишется в виде


\begin{equation}


Aa=\int_\Gamma g\alpha.


\end{equation}


В силу единственности решения задачи Римана (1) ранг $r$ матрицы $A$


должен быть равен числу неизвестных $a_{kn}$, т.\,е.


$r=\kappa_1+\kappa_2+2\ord\Delta +2$.





Пусть $B$ --- невырожденная квадратная матрица порядка $r$,


составленная из строк матрицы $A$ с номерами $j_1, j_2,\dotsc, j_r$


($j_1<j_2<\dotsb <j_r$), $B_j$ --- квадратная матрица порядка $r+1$,


составленная из $r+1$ строк расширенной матрицы $(A,\ \int\limits_\Gamma


g\alpha)$ с номерами $j_1, j_2,\dotsc,j_r, j$. Тогда


$$


\det B_j=\int_\Gamma g\beta_j,


$$


где


$$


\beta_j=\sum_{n=1}^{r}B_{j,j_n}\alpha_{j_n} +B_{j,j}\alpha_j,


$$


$B_{j,k}$ --- алгебраическое дополнение элемента


$\int\limits_\Gamma g\alpha_k$ матрицы


$B_j$; очевидно, $B_{j,j}=\det B\ne 0$. Если $j$ принимает одно из


значений $j_1, j_2,\dotsc,j_r$, то $\beta_j=0$. Условия разрешимости


системы (11) (или (12)) имеют вид


\begin{equation}


\int_\Gamma g\beta_j=0,\quad j\in \bold Z,


\ \; j\ge 0,\ \; j\ne j_1,j_2,\dotsc, j_r.


\end{equation}


Совокупность условий (10) и (13) необходима и достаточна для


разрешимости задачи (1). При их выполнении единственное решение задачи


(1) определяется равенством $F(q)=f_k(z(q))$, где функция $f_1=f_2$


определена формулой (8). Таким образом, доказана





\begin{thm}


Пусть $\gamma$ имеет положительную линейную меру и


коэффициент $G$ задачи $(1)$ удовлетворяет неравенству


$G(j(t))\not\equiv G(t)$, $t\in \Gamma_1\cup\Gamma_2$. Тогда для


разрешимости задачи Римана $(1)$ необходимы и достаточны


условия $(10)$ и $(13)$. При их выполнении задача $(1)$ имеет единственное


решение, определяемое равенством $F(q)=f_k(z(q))$, где совпадающие


между собой функции $f_k$, $k=1,\,2$, определяются равенством $(8)$.


\end{thm}





\section{Случай произвольного кусочно-гладкого контура}





1. Пусть теперь $\Gamma$ --- произвольный кусочно-гладкий контур на $R$.


Зададим дивизор $D$, носитель которого лежит в $R\setminus \Gamma$, и


функции $G \in H_{\mu, (\mu)}(\Gamma, T)$ и


$g \in H_{\mu,\lambda}(\Gamma, T)$, $0<\mu < 1$, $-1<\lambda <0$,


причем $G(t)\ne 0$, $t\in \Gamma$. Рассмотрим краевую задачу Римана


в следующей постановке.





{\it


Найти кусочно-мероморфную функцию $F$ на $R$ с линией скачков


$\Gamma$, кратную дивизору $1/D$ и ограниченную в окрестности


идеальной границы поверхности $R$, предельные значения которой на


$\Gamma$ принадлежат классу $H_{\mu, \lambda}(\Gamma, T)$,


$0<\mu<1$, $-1<\lambda <0$, и удовлетворяют соотношению


}


\begin{equation}


F^+(t)=G(t)F^-(t)+g(t),\quad t\in \Gamma.


\end{equation}





2. Обозначим через $E$ относительно компактное открытое множество


на $R$, обладающее следующими свойствами: 1)~граница $\partial E$


множества $E$ есть аналитическая жорданова кривая, состоящая из двух


компонент, гомеоморфных окружности;


2)~$\partial E$ не содержит точек ветвления накрывающей


$(R, z)$; 3)~$j_R(E)=E$; 4)~$\Gamma \subset E$; 5)~$E$ содержит четное


число точек ветвления накрывающей $(R, z)$; 6)~среди множеств,


обладающих свойствами 1)--5), $E$ содержит минимальное число точек


ветвления накрывающей $(R, z)$.





Обозначим через $\alpha_k$, $k=1,\ 2,\dotsc, 2h+2$, точки


ветвления накрывающей $(R, z)$, лежащие в $E$. В случае $h=-1$\; $E$


состоит из двух односвязных компонент, каждая из которых однолистно


и конформно отображается в $\Bbb C$ функцией $z: E \to \Bbb C$.


В случае $h\ge 0$\; $E$ есть область рода $h$.





3. Рассмотрим сначала случай, когда $h\ge 0$. Через $(S, z)$


обозначим двулистную накрывающую замкнутой комплексной плоскости


$\ov{\Bbb C}$, определяемую уравнением


\begin{equation}


w^2=\prod_{k=1}^{2h+2}(z-r_k),


\end{equation}


где $r_k=z(\alpha_k)$.\; $S$ есть гиперэллиптическая риманова


поверхность рода $h$. Через $E_0$ обозначим область на $S$, являющуюся


полным прообразом области $z(E)\subset \Bbb C$ относительно


отображения $z: S\to \ov{\Bbb C}$. Ясно, что $E_0$ содержит


все точки ветвления накрывающей $(S, z)$ и существует конформный


гомеоморфизм $\alpha :E_0\to E$ такой, что $z\circ \alpha =z$ и


$\alpha\circ j_S=j_R\circ \alpha$. $S\setminus \ov E_0$


состоит из двух односвязных компонент $E_1$ и $E_2$. Отображение


$z: E_k\to \ov{\Bbb C}$, $k=1,\, 2$, однолистно и конформно;


обратное к нему отображение обозначим через $\rho_k: z(E_k)\to E_k$.


Очевидно, $\rho_2=j_S\circ \rho_1$.





Положим $\Gamma_0=\alpha^{-1}(\Gamma)$ и определим на $\Gamma_0$


ориентацию, индуцированную отображением $\alpha^{-1}: \Gamma\to


\Gamma_0$. Пусть $F$ --- решение задачи Римана (14). Легко видеть,


что $F\circ j_R=F$ в $R\setminus E$. Поэтому на $S$ определена


однозначная кусочно-мероморфная функция


$$


F_0(q)=


\begin{cases}


F(\alpha (q)), & q\in E_0;\\


F(p), & z(p)=z(q),\ q\in S\setminus E_0,\ p\in R\setminus E,


\end{cases}


$$


с линией скачков $\Gamma_0$. Определим на $S$ дивизор $D_0$, полагая


$$


\ord_qD_0=


\begin{cases}


\ord_{\alpha (q)}D, & q\in E_0;\\


\min(\ord_pD,\, \ord_{j_R(p)}D), & z(p)=z(q),


\ q\in S\setminus E_0,\ p\in R\setminus E,\ j_R(p)\ne p;\\


\lbrack \frac{1}{2}


\ord_p D], & z(p)=z(q),\ q\in S\setminus E_0,\ p\in


R\setminus E,\ j_R(p)=p.


\end{cases}


$$





Пусть $T_0=\alpha^{-1}(T)$ и доопределим на $T_0$ функцию


$\lambda$, полагая $\lambda |T_0=\lambda\circ\alpha |T_0$. Функция


$F_0$ мероморфна в $S\setminus \Gamma_0$, кратна дивизору $1/D_0$, а


ее предельные значения на $\Gamma_0$ принадлежат классу


$H_{\mu,\lambda}(\Gamma_0,T_0)$ и удовлетворяют соотношению вида


\begin{equation}


F_0^+(t)=G(\alpha(t))F_0^-(t)+g(\alpha(t)),\quad t\in \Gamma_0.


\end{equation}


Кроме того, в $S\setminus E_0$ она удовлетворяет соотношению


\begin{equation}


F_0\circ j_S=F_0.


\end{equation}


Функции $G\circ\alpha$ и $g\circ\alpha$ принадлежат классам


$H_{\mu, (\mu)}(\Gamma_0, T_0)$ и $H_{\mu, \lambda}(\Gamma_0, T_0)$


соответственно.





Существует взаимнооднозначное соответствие между точками $q$


поверхности $S$ и парами чисел $(z, w)=(z(q), w(z(q)))$, связанными


соотношением (15). Точку $q$ обычно отождествляют с парой $(z, w)$.


Тогда точке $j_S(q)$ соответствует пара $(z, -w)$. Точке ветвления


$\alpha_k$ соответствует пара $(r_k, 0)$.





Выберем на $S$ канонические циклы $\{a_k,\, b_k\}$,


$k=1,2,\dotsc,h$, как в [2]. Тогда комплексно нормированный базис


пространства абелевых дифференциалов первого рода на $S$ имеет вид [2]


\begin{gather*}


\varphi_j(q)=


\frac{


\begin{vmatrix}


0 & \frac{dz}w & \frac{zdz}w & \dotsb & \frac{z^{h-1}dz}w \\[2mm]


-\delta_{j1} & \int \limits_{a_1}\frac{d\tau}\zeta &


\int \limits_{a_1}\frac{\tau d\tau}\zeta & \dotsb &


\int \limits_{a_1}\frac{\tau^{h-1}d\tau}\zeta \\


\vdots & \vdots & \vdots & \ddots & \vdots \\


-\delta_{jh} & \int \limits_{a_h}\frac{d\tau}\zeta &


\int \limits_{a_h}\frac{\tau d\tau}\zeta & \dotsb &


\int \limits_{a_h}\frac{\tau^{h-1}d\tau}\zeta


\end{vmatrix}


}


{\begin{vmatrix}


\int \limits_{a_1}\frac{d\tau}\zeta &


\int \limits_{a_1}\frac{\tau d\tau}\zeta & \dotsb &


\int \limits_{a_1}\frac{\tau^{h-1}d\tau}\zeta \\


\vdots & \vdots & \ddots & \vdots \\


\int \limits_{a_h}\frac{d\tau}\zeta &


\int \limits_{a_h}\frac{\tau d\tau}\zeta & \dotsb &


\int \limits_{a_h}\frac{\tau^{h-1}d\tau}\zeta


\end{vmatrix}


}, \\


%


j=1, 2,\dotsc,h,\ \; q=(z, w),\ \; t=(\tau, \zeta)\in S.


\end{gather*}


Нормированный абелев интеграл третьего рода, служащий разрывным


аналогом ядра Коши, имеет вид [2]


$$


\omega_{qq_0}(t)=\omega_q(t)-\omega_{q_0}(t),


$$


где


$$


\omega_q(t)=


\frac{


\begin{vmatrix}


\frac{(w+\zeta)d\tau}{2\zeta (\tau-z)} &


\frac{d\tau}\zeta & \frac{\tau d\tau}\zeta & \dotsb &


\frac{\tau^{h-1}d\tau}\zeta \\[2mm]


\int \limits_{a_1}\frac{(w+\zeta)d\tau}{2\zeta (\tau-z)} &


\int \limits_{a_1}\frac{d\tau}\zeta &


\int \limits_{a_1}\frac{\tau d\tau}\zeta & \dotsb &


\int \limits_{a_1}\frac{\tau^{h-1}d\tau}\zeta \\


\vdots & \vdots & \vdots & \ddots & \vdots \\


\int \limits_{a_h}\frac{(w+\zeta)d\tau}{2\zeta (\tau-z)} &


\int \limits_{a_h}\frac{d\tau}\zeta &


\int \limits_{a_h}\frac{\tau d\tau}\zeta & \dotsb &


\int \limits_{a_h}\frac{\tau^{h-1}d\tau}\zeta


\end{vmatrix}


}


{


\begin{vmatrix}


\int \limits_{a_1}\frac{d\tau}\zeta &


\int \limits_{a_1}\frac{\tau d\tau}\zeta & \dotsb &


\int \limits_{a_1}\frac{\tau^{h-1}d\tau}\zeta \\


\vdots & \vdots & \ddots & \vdots \\


\int \limits_{a_h}\frac{d\tau}\zeta &


\int \limits_{a_h}\frac{\tau d\tau}\zeta & \dotsb &


\int \limits_{a_h}\frac{\tau^{h-1}d\tau}\zeta


\end{vmatrix}


}.


$$





Положим


$$


X(q)=\exp\bigg(\frac1{2\pi i}\int_{\Gamma_0}\ln G(\alpha


(t))\omega_{qq_0}(t)\bigg),\quad q\in S.


$$


Нули и бесконечности функции $X(q)$ образуют квазидивизор


$(X)=\tau_1^{\kappa_1}\tau_2^{\kappa_2}\dotsb \tau_r^{\kappa_r}$,


где $\tau_k=\alpha^{-1}(t_k)$, $t_k\in T$, --- узлы линии $\Gamma_0$,


$\kappa_k$ --- числа, определяемые коэффициентом $G$ и выбором ветви


функции $\ln G(\alpha (t))$ [2]. Положим


$\kappa=\sum\limits_{k=1}^r[\kappa_k-\lambda(\tau_k)]$, где $[\ ]$ означает


целую часть числа. Число $\kappa$, не


зависящее от выбора ветви $\ln G\circ\alpha$, назовем индексом


коэффициента задачи (16) в классе $H_{\mu,\lambda}(\Gamma_0, T_0)$.


Положим


$A=t_1^{[\kappa_1-\lambda(\tau_1)]}t_2^{[\kappa_2-\lambda(\tau_2)]}\dotsb


t_r^{[\kappa_r-\lambda(\tau_r)]}$, $B=(q')^hq_1^{-1}\dotsb q_h^{-1}$,


где $q'\in S$ --- произвольно фиксированная точка, не совпадающая с


$q_0$, точки $q_1,q_2,\dotsc,q_h$ образуют решение проблемы обращения


Якоби вида


$$


\sum_{j=1}^h\int_{q'}^{q_j}\varphi_\nu\equiv \frac1{2\pi i}


\int_{\Gamma_0}\ln G(\alpha (t))\varphi_\nu (t)


\ \; (\text{по модулю периодов}),\ \; \nu=1,\ 2,\dotsc,h.


$$


Тогда общее решение однородной $(g=0)$ задачи Римана (16) имеет вид


\begin{equation}


f(q)\exp\bigg(\frac1{2\pi i}\int_{\Gamma_0}\ln G(\alpha


(t))\omega_{qq_0}(t)


-\sum_{j=1}^h\Big(\int_{q'}^{q_j}\omega_{qq_0}


-2\pi im_j\int_{q'}^{q}\varphi_j\Big)\bigg),


\end{equation}


где в последних двух интегралах путь интегрирования не


пересекает канонических сечений $a_1,a_2,\dotsc,a_h$, $m_j$ --- вполне


определенные целые числа, $f$ --- произвольная мероморфная функция,


кратная дивизору $D_0^{-1}A^{-1}B^{-1}$. У функции (по переменной


$q$) $\omega_{qq_0}$ выбрана фиксированная в $S\setminus


\mathop{\cup}\limits_{k=1}^ha_k$ ветвь, исчезающая в точке $q_0$.





Обозначим через $l_0$ и $l'_0$ число линейно независимых


мероморфных функций и дифференциалов на $S$, кратных соответственно


дивизорам $D_0^{-1}A^{-1}B^{-1}$ и $D_0AB$. Через $P$ обозначим целый


дивизор порядка $l'_0$ с носителем в $S\setminus\Gamma_0$ и такой, что


не существует абелевых дифференциалов на $S$, кратных дивизору


$D_0ABP$ и отличных от тождественного нуля. Для построения решения


неоднородной задачи (16) найдем сначала частное решение этой задачи в


классе функций, кратных дивизору $D_0^{-1}P^{-1}$. Такая задача


безусловно разрешима в силу выбора дивизора $P$. Ее частным решением


является функция вида [2]


\begin{equation}


\frac{X_0(q)}{2\pi i}


\int_{\Gamma_0}\frac{g(\alpha(t))}{X_0^+(t)}A_1(t, q),


\end{equation}


где $X_0$ --- функция вида (18) при $f=1$, $A_1(t, q)$ ---


мероморфный аналог ядра Коши на $S$ с характеристическим дивизором


$\Delta_1$, который получается делением дивизора $D_0ABP$ на некоторый


целый дивизор. Функция (19) будет решением задачи (16) в том и только


том случае, если $g$ удовлетворяет $l_0'$ условиям, обеспечивающим


кратность функции (19) дивизору $D^{-1}$. Эти условия равносильны


условиям разрешимости задачи (16) и имеют вид


\begin{equation}


\int_{\Gamma_0}g\circ\alpha\frac{\psi_j}{X_0}=0,\quad


j=1,2,\dotsc,l_0',


\end{equation}


где $\psi_j$, $j=1, 2,\dotsc,l'_0$, --- базис пространства абелевых


дифференциалов на $S$, кратных дивизору $D_0AB$.





При выполнении условия (20) общее решение задачи (16) имеет вид


\begin{multline}


F_0(q)=


f(q)\exp\bigg(\frac1{2\pi i}\int_{\Gamma_0}\ln G(\alpha


(t))\omega_{qq_0}(t)-\sum_{j=1}^h\Big(\int_{q'}^{q_j}\omega_{qq_0}


-2\pi im_j\int_{q'}^{q}\varphi_j\Big)\bigg)+ \\


%


+\frac{X_0(q)}{2\pi i}


\int_{\Gamma_0}\frac{g(\alpha(t))}{X_0^+(t)}A_1(t, q),\quad q\in S.


\end{multline}


Чтобы эта функция определяла решение задачи (14), необходимо и


достаточно условия (17).





Обозначим через $f_1,f_2,\dotsc,f_{l_0}$ базис пространства


мероморфных функций на $S$, кратных дивизору $D_0^{-1}A^{-1}B^{-1}$.


Тогда функция $F_0$ может быть записана в виде


\begin{equation}


F_0(q)=X_0(q)\sum_{k=1}^{l_0}c_kf_k(q)+


\frac{X_0(q)}{2\pi i}


\int_{\Gamma_0}\frac{g(\alpha(t))}{X_0^+(t)}A_1(t, q),


\end{equation}


где $c_k$ --- произвольные комплексные числа.





Пусть $s$ есть фиксированная точка в $S\setminus E_0$


такая, что $z(s)\ne \infty$ и функции $X_0(q)$, $f_k(q)$,


$A_1(t, q)$ голоморфны по $q$ в точках $s$ и $j_S(s)$.


Используем ряды Тейлора


$$


X_0(q)f_k(q)=\sum_{i=0}^\infty a_{ik}(z(q)-z(s))^i,\quad


X_0(j_S(q))f_k(j_S(q))=\sum_{i=0}^\infty b_{ik}(z(q)-z(s))^i.


$$


Аналогично в окрестности точки $q=s$ имеем


$$


\frac{X_0(q)A_1(t, q)}{2\pi iX_0^+(t)}=\sum_{i=0}^\infty


\gamma_i(t)(z(q)-z(s))^i,\quad


\frac{X_0(j_S(q))A_1(t, j_S(q))}{2\pi iX_0^+(t)}=\sum_{i=0}^\infty


\delta_i(t)(z(q)-z(s))^i.


$$





В силу этих разложений равенство (17) равносильно


системе линейных алгебраических уравнений


\begin{equation}


\sum_{k=1}^{l_0}(a_{ik}-b_{ik})c_k=-\int_{\Gamma_0}g(\alpha(t))(\gamma_i(t)-


\delta_i(t)),\quad i=0,1,2,\dotsc


\end{equation}


Полагая $U=a_{ik}-b_{ik}$, $i=0,1,2,\dotsc$, $k=1,2,\dotsc,l_0$,


$\beta_i=-(\gamma_i-\delta_i)\circ \alpha^{-1}$,


$\beta=(\beta_0,\beta_1,\dotsc)^t$,


$c=(c_1,c_2,\dotsc,c_{l_0})^t$,


перепишем систему


(23) в виде


\begin{equation}


Uc=\int_\Gamma g\beta.


\end{equation}





Обозначим через $r$ ранг матрицы $U$. Пусть $V$ --- невырожденная


квадратная матрица порядка $r$, составленная из элементов матрицы $U$,


стоящих на пересечении строк с номерами $i_1,i_2,\dotsc,i_r$


($i_1<i_2<\dotsb <i_r$) и столбцов с номерами $k_1,k_2,\dotsc,k_r$


($k_1<k_2<\dotsb <k_r$), $V_i$ --- квадратная матрица порядка $r+1$,


состоящая из элементов расширенной матрицы $(U,\int\limits_\Gamma g\beta)$,


стоящих на пересечении строк с номерами $i_1,i_2,\dotsc,i_r,\,i$ и


столбцов с номерами $k_1,k_2,\dotsc,k_r,\,l_0+1$. Тогда


$\det V_i=\int\limits_\Gamma g\lambda_i$, где


$$


\lambda_i=\sum_{n=1}^rV_{i,i_n}\beta_{i_n}+V_{i,i}\beta_i,


$$


$V_{i,j}$ --- алгебраическое дополнение элемента


$\int\limits_\Gamma g\beta_j$


матрицы $V_i$, $V_{i,i}=\det V$. Условия разрешимости системы (24)


имеют вид


\begin{equation}


\int_\Gamma g\lambda_i=0,\ \; i\in \Bbb Z,\ \; i\ge 0,


\ \; i\ne i_1,i_2,\dotsc,i_r.


\end{equation}


При выполнении условий (25) числа $c_{k_1}, c_{k_2},\dotsc,c_{k_r}$


выражаются линейно через $\int\limits_\Gamma g\beta_i$,


$i=i_1,i_2,\dotsc,i_r$, т.\,е. являются ограниченными линейными функционалами


в пространстве $H_{\mu,\lambda}(\Gamma,T)$, а остальные числа $c_k$


остаются произвольными.





Условия (20) перепишем в виде


\begin{equation}


\int_{\Gamma}g\frac{\psi_j\circ \alpha^{-1}}{X_0\circ \alpha^{-1}}=0,


\quad j=1, 2,\dotsc,l_0'.


\end{equation}


Если функция $g$ удовлетворяет условиям (25) и (26), то функция $F_0$


из соотношения (22) определяет решение $F$ задачи (14) по


формуле


\begin{equation}


F(q)=\begin{cases}


F_0(\alpha^{-1}(q)), & q\in E;\\


F_0(p), & z(p)=z(q),\ q\in R\setminus E,\ p\in S\setminus E_0.


\end{cases}


\end{equation}


Полученные результаты формулируются как





\begin{thm}


Пусть $h\ge 0$. Тогда для разрешимости задачи Римана


$(14)$ необходимы и


достаточны условия $(25)$ и $(26)$. При их выполнении решение


задачи определяется равенствами $(22)$ и $(27)$, причем в равенстве


$(22)$\; $r$


из чисел $c_k$ являются линейными ограниченными функционалами от $g\in


H_{\mu,\lambda}(\Gamma,T)$, а остальные произвольны. Однородная $(g=0)$


задача Римана $(14)$ имеет $l_0-r$ линейно независимых решений.


\end{thm}





4. Рассмотрим теперь случай, когда $h=-1$, т.\,е. $E$ состоит из двух


компонент $O_1$ и $O_2=j_R(O_1)$, гомеоморфных открытому кругу, $S$


есть дизъюнктное объединение двух экземпляров замкнутой комплексной


плоскости $\ov{\Bbb C}$. Положим


$\Gamma_k=\Gamma\cap O_k$, $\Gamma_{0k}=z(\Gamma_k)$, $k=1,\,2$. Тогда


задача (16) распадается на две задачи Римана на $\ov{\Bbb C}$


с краевыми


условиями на контурах $\Gamma_{0k}$. Отображение $z: O_k\to\ov{\Bbb C}$


однолистно и конформно. Обратное к нему отображение обозначим через


$\rho_k$. Краевые условия соответствующих задач будут иметь вид


\begin{equation}


F_{0k}^+(t)=G(\rho_k(t))F^-_{0k}(t)+g(\rho_k(t)),\ \; t\in\Gamma_{0k},


\ \; k=1,\,2. \tag{28$_k$}


\end{equation}


\addtocounter{equation}{1}





Решение задачи $(28_k)$ ищется в классе кусочно-мероморфных функций,


кратных дивизору $1/D_{0k}$, где


$$


\ord_{z(q)}D_{0k}=\begin{cases}


\ord_qD, & q\in O_k;\\


\min(\ord_qD,\, \ord_{j_R(q)}D), & q\in R\setminus E,\ j_R(q)\ne q;\\


\lbrack\frac12\ord_qD], & q\in R\setminus E,\ j_R(q)=q.


\end{cases}


$$


На $\Gamma_{0k}$ решение задачи $(15_k)$ должно принадлежать классу


$H_{\mu,\lambda}(\Gamma_{0k},T_{0k})$, где $T_{0k}=z(T\cap O_k)$, $k=1,\,2$.


Предположим для определенности, что $T\cap O_1=\{t_1,\dotsc,t_m\}$,


$T\cap O_2=\{t_{m+1},\dotsc,t_r\}$. Чтобы функции $F_{01}$ и $F_{02}$


определяли решение задачи (14), они должны удовлетворять условию


\begin{equation}


F_{01}(z)=F_{02}(z),\quad z\in \ov{\Bbb C}\setminus z(E).


\end{equation}





Положим


\begin{gather*}


X_k(z)=\exp\bigg(\frac1{2\pi i}\int_{\Gamma_{0k}}\frac{\ln


G(\rho_k(\tau))d\tau}{\tau-z}\bigg), \\


%


\kappa_{ki}=\ord_{z(t_i)}X_k\ \; (i=1,\dotsc,m


\text{\; при \;} k=1 \text{\; и \;}


i=m+1,\dotsc,r\text{\; при \;} k=2), \\


%


Y_1=\prod_{i=1}^m(z-z(t_i))^{[\kappa_{1i}-\lambda(t_i)]},\quad


Y_2=\prod_{i=m+1}^r(z-z(t_i))^{[\kappa_{2i}-\lambda(t_i)]}, \\


X_{0k}=X_k/Y_k,\ \; k=1,\, 2.


\end{gather*}


Функция $X_{0k}$ удовлетворяет однородному краевому условию $(28_k)$, а ее


нули и бесконечности в $\Bbb C$ образуют квазидивизор


$$


\prod_{i=1}^m(z(t_i))^{\kappa_{1i}-[\kappa_{1i}-\lambda(t_i)]}


\text{\; при \;} k=1 \quad\text{и}\quad


\prod_{i=m+1}^r(z(t_i))^{\kappa_{2i}-[\kappa_{2i}-\lambda(t_i)]}


\text{\; при \;} k=2.


$$


Общее решение однородной задачи $(28_k)$ имеет вид $X_{0k}P_k$,


где $P_k$ --- произвольный многочлен степени не выше, чем


$\kappa^{(k)}=-\ord_\infty Y_k$


(многочлен отрицательной степени считаем равным нулю).





Решение неоднородной задачи $(28_k)$ при $\kappa^{(k)}\ge-1$ существует


при любой функции $g\in H_{\mu,\lambda}(\Gamma, T)$ и может быть записано в


виде


\begin{equation}


F_{0k}(z)=\frac{X_{0k}(z)}{2\pi i}


\int_{\Gamma_{0k}}\frac{g(\rho_k(\tau))d\tau}{X_{0k}^+(\tau)(\tau-z)}+


X_{0k}(z)P_k(z),\quad k=1,\, 2.


\end{equation}


При $\kappa^{(k)}<-1$ для разрешимости задачи $(28_k)$ необходимо и


достаточно, чтобы функция $g$ удовлетворяла условиям


\begin{equation}


\int_{\Gamma_{0k}}\frac{\tau^ig(\rho_k(\tau))d\tau}{X_{0k}^+(\tau)}=0,


\quad i=0,\ 1,\dotsc,-\kappa^{(k)}-2.


\end{equation}


При их выполнении задача $(28_k)$ имеет единственное решение, определяемое


формулой (30).





Чтобы функции (30) определяли решение задачи (14), они должны


удовлетворять соотношению (29). Для выполнения этих условий разложим


функцию $F_{0k}$ в ряд Тейлора в окрестности некоторой точки $a\in


\Bbb C\setminus z(E)$ голоморфности $X_{0k}$.





Запишем ряд Тейлора


$$


\frac{X_{0k}(z)}{2\pi iX^+_{0k}(\tau)(\tau -z)}=


\sum_{j=0}^\infty c_{kj}(\tau)(z-a)^j


$$


и положим


$$


P_k(z)=\sum_{n=0}^{\kappa^{(k)}}a_{kn}z^n


$$


(если верхний индекс суммирования меньше нижнего, то сумма считается


пустой). Разлагая функции $X_{0k}(z)z^n$ в ряды Тейлора, получим


$$


X_{0k}(z)z^n=\sum_{j=0}^{\infty}a_{kjn}(z-a)^j.


$$


Сравнивая коэффициенты тейлоровских разложений функций $F_{01}$ и $F_{02}$,


получим соотношения, эквивалентные равенству (29),


\begin{equation}


\sum_{n=0}^{\kappa^{(1)}}a_{1jn}a_{1n}-


\sum_{n=0}^{\kappa^{(2)}}a_{2jn}a_{2n}=


-\int_{\Gamma_{01}}g(\rho_1(\tau ))c_{1j}(\tau )\d\tau +


\int_{\Gamma_{02}}g(\rho_2(\tau ))c_{2j}(\tau )\d\tau,\ \;j=0,1,2,\dotsc


\end{equation}


Положим


$$


A=


\begin{pmatrix}


a_{100} & a_{101} & \dotsb & a_{10\kappa^{(1)}} &


-a_{200} & -a_{201} & \dotsb & -a_{20\kappa^{(2)}} \\


a_{110} & a_{111} & \dotsb & a_{11\kappa^{(1)}} &


-a_{210} & -a_{211} & \dotsb & -a_{21\kappa^{(2)}} \\


a_{120} & a_{121} & \dotsb & a_{12\kappa^{(1)}} &


-a_{220} & -a_{221} & \dotsb & -a_{22\kappa^{(2)}} \\


\vdots & \vdots & \ddots & \vdots & \vdots & \vdots & \ddots & \vdots


\end{pmatrix}


$$


(если $\kappa^{(k)}<0$, то столбцы, содержащие $a_{kij}$, в матрице


$A$ отсутствуют),


\begin{gather*}


\alpha_j(t)=\begin{cases}


-c_{1j}(z(t))dz(t), & t\in\Gamma_1;\\


\phantom{-}c_{2j}(z(t))dz(t), & t\in\Gamma_2,


\end{cases}


\\


\alpha (t)=(\alpha_1(t),\ \alpha_2(t),\dotsc)^t, \\


%


a=(a_{10},\, a_{11},\dotsc,a_{1\, \kappa^{(1)}},


\, a_{20},\, a_{21},\dotsc,a_{2\, \kappa^{(2)}})^t.


\end{gather*}


Тогда система (32) запишется в виде


\begin{equation}


Aa=\int_\Gamma g\alpha.


\end{equation}


Обозначим через $r$ ранг матрицы $A$. Пусть $B$ --- невырожденная


квадратная матрица порядка $r$, составленная из элементов матрицы $A$,


стоящих на пересечении строк с номерами $j_1,j_2,\dotsc,j_r$


($j_1<j_2<\dotsb<j_r$) и столбцов с номерами $k_1,k_2,\dotsc,k_r$ ($k_1<


k_2<\dotsb<k_r$), $B_j$ --- квадратная матрица порядка $r+1$, состоящая


из элементов расширенной матрицы $(A,\int\limits_\Gamma g\beta)$, стоящих на


пересечении строк с номерами $j_1,j_2,\dotsc,j_r,\,j$ и столбцов с


номерами $k_1,k_2,\dotsc,k_r,\, \kappa_+^{(1)}+\kappa_+^{(2)}+3$, где


$\kappa_+^{(k)}:=\max\{-1,\kappa^{(k)}\}$. Тогда


$$


\det B_j=\int_\Gamma g\beta_j,


\quad\text{где}\quad


\beta_j=\sum_{n=1}^{r}B_{j,j_n}\alpha_{j_n} +B_{j,j}\alpha_j,


$$


$B_{j,k}$ --- алгебраическое дополнение элемента


$\int\limits_\Gamma g\alpha_k$


матрицы $B_j$; очевидно, $B_{j,j}=\det B\ne 0$. Если $j$ принимает одно из


значений $j_1, j_2,\dotsc,j_r$, то $\beta_j=0$. Условия разрешимости


системы (32) (или (33)) имеют вид


\begin{equation}


\int_\Gamma g\beta_j=0,\quad


j\in \Bbb Z,\ \; j\ge 0,\ \; j\ne j_1,j_2,\dotsc, j_r.


\end{equation}





Условия (31) перепишем в виде


\begin{equation}


\int_{\Gamma_k}\frac{(z(t))^ig(t)dz(t)}{X^+_{0k}(z(t))}=0,\quad


i=0,1,\dotsc,-\kappa^{(k)}-2,\quad k=1,\,2.


\end{equation}





Совокупность условий (34) и (35) необходима и достаточна для


разрешимости задачи (14). При их выполнении общее решение задачи


(14) определяется равенством


\begin{equation}


F(q)=


\begin{cases}


F_{01}(z(q))=F_{02}(z(q)), & q\in R\setminus E;\\


F_{0k}(z(q)), & q\in O_k,\ \; k=1,\,2,


\end{cases}


\end{equation}


где $F_{0k}$ определена формулой (30).





Таким образом, доказана





\begin{thm}


Пусть $h=-1$. Тогда для


разрешимости задачи Римана $(14)$ необходимы и достаточны


условия $(34)$ и $(35)$. При их выполнении общее решение задачи $(14)$


определяется формулами $(30)$ и $(36)$, где коэффициенты $a_{kn}$


многочленов $P_k$ удовлетворяют системе линейных алгебраических


уравнений $(32)$. Число линейно независимых решений однородной задачи


Римана $(14)$ равно $\kappa_+^{(1)}+\kappa_+^{(2)}-r+2$.


\end{thm}





\begin{thebibliography}{9}





\bibitem{1}


Чибрикова Л.И. {\em Граничные задачи теории аналитических функций


на римановых поверхностях} // Итоги науки и техн. ВИНИТИ.


Матем. анализ. -- 1980. -- Т.\,18. -- С.\,3--66.


\bibitem{2}


Зверович Э.И. {\em Краевые задачи теории аналитических функций в


г\"ельдеровских классах на римановых поверхностях } // УМН. -- 1971.


-- Т.\,26. -- Вып.\,1. -- С.\,113--179.


\bibitem{3}


Бикчантаев И.А. {\em Аналоги ядра Коши на римановой поверхности и


некоторые их приложения} // Матем. сб. -- 1980. --


Т.\,112. -- \No\,2. -- С.\,256--282.


\bibitem{4}


Бикчантаев И.А. {\em О числе решений краевой задачи Римана на


некомпактной римановой поверхности} // Изв. вузов. Математика.


-- 1984. -- \No\,9. -- С.\,10--13.


\bibitem{5}


Бикчантаев И.А. {\em Обращение сингулярных интегралов на


кусочно-гладком контуре на римановой поверхности} // Дифференц.


уравнения. -- 1994. -- Т.\,30. -- \No\,5. -- С.\,882--888.


\bibitem{6}


Солдатов А.П. {\em Одномерные сингулярные операторы и краевые


задачи теории функций}. -- М.: Высш. школа, 1991. -- 208 c.


\bibitem{7}


Sario L., Nakai M. {\em Classification theory of open Riemann


surfaces}. -- Berlin--Heidelberg--N.~Y.: Springer--Verlag, 1970. -- 446 p.





\end{thebibliography}





\vskip 2cm





\postup{\it Казанский государственный}{$\qquad$\it Поступили}


\postup{\it университет}{{\it первый вариант} 11.04.1996}


\postup{}{{\it окончательный вариант} 05.05.1998}





\label{end}


\end{document}


